\subsection{\texorpdfstring{DERIVATIVES OF THE CIRCULAR FUNCTIONS; THE NUMBER \(\pi\)}{DERIVATIVES OF THE CIRCULAR FUNCTIONS; THE NUMBER PI}}

We have defined, in General Topology (Gen.~Top., VIII, p.~106), the continuous homomorphism \(x \mapsto \mathbf{e}(x)\) of the additive group \(\mathbf{R}\) onto the multiplicative group \(\mathbf{U}\) of complex numbers of absolute value 1; this is a periodic function with principal period 1, and \(\mathbf{e}\left(\frac{1}{4}\right) = i\) . One knows (loc. cit.) that every continuous homomorphism of \(\mathbf{R}\) onto \(\mathbf{U}\) is of the form \(x \mapsto \mathbf{e}(x / a)\) , and one puts \(\cos_a x = \mathcal{R}(\mathbf{e}(x / a))\) , \(\sin_a x = \mathcal{I}(\mathbf{e}(x / a))\) (trigonometric functions, or circular functions, to base \(a\) ); these last functions are continuous maps from \(\mathbf{R}\) into \([-1, +1]\) having principal period \(a\) . We have \(\sin_a(x + a / 4) = \cos_a x\) , \(\cos_a(x + a / 4) = -\sin_a x\) , and the function \(\sin_a x\) is increasing on the interval \([-a / 4, a / 4]\) .

PROPOSITION 3. The function \(\mathbf{e}(x)\) has a derivative equal to \(2\pi i\mathbf{e}(x)\) at every point of R, where \(\pi\) is a constant \textgreater0.

Now, th. 1 of III, p.~51, applied to the case where E is the field C of complex numbers, yields the relation \(\mathbf{e}'(x) = \mathbf{e}'(0)\mathbf{e}(x)\) ; moreover, since \(\mathbf{e}(x)\) has constant euclidean norm, \(\mathbf{e}'(x)\) is orthogonal to \(\mathbf{e}(x)\) (I, p.~7, Example 3); one thus has \(\mathbf{e}'(0) = \alpha i\) , with \(\alpha\) real. Since \(\sin_1 x\) is increasing on \([- \frac{1}{4}, \frac{1}{4}]\) its derivative for \(x = 0\) is \(\geqslant 0\) , so \(\alpha \geqslant 0\) , and since \(\mathbf{e}(x)\) is not constant, \(\alpha > 0\) ; it is conventional to denote the number \(\alpha\) so obtained by \(2\pi\) .

In §2 (III, p.~23) we shall show how one can calculate arbitrarily close approximations to the number \(\pi\) .

We thus have the formula

\[
\mathrm{D} \left(\mathbf {e} \left(\frac {x}{a}\right)\right) = \frac {2 \pi i}{a} \mathbf {e} \left(\frac {x}{a}\right).\tag{12}
\]

One sees that this formula simplifies when \(a = 2\pi\) ; this is why one uses the circular functions relative to base \(2\pi\) exclusively in Analysis; we agree to omit the base in the notation for these functions; unless mentioned expressly to the contrary, the notations \(\cos x\) , \(\sin x\) and \(\tan x\) denote \(\cos_{2\pi}x\) , \(\sin_{2\pi}x\) and \(\tan_{2\pi}x\) respectively.

With these conventions, and \(a = 2\pi\) , formula (12) can be written

\[
\mathrm{D} (\cos x + i \sin x) = \cos \left(x + \frac {\pi}{2}\right) + i \sin \left(x + \frac {\pi}{2}\right),\tag{13}
\]

which is equivalent to

\[
\mathrm{D} (\cos x) = - \sin x, \quad \mathrm{D} (\sin x) = \cos x,
\]

from which one deduces

\[
\mathrm{D} (\tan x) = 1 + \tan^ {2} x = \frac {1}{\cos^ {2} x}.\tag{15}
\]

Besides the three circular functions \(\cos x\) , \(\sin x\) and \(\tan x\) one also uses, in numerical work, the three auxiliary functions: cotangent, secant and cosecant, defined by the formulae

\[
\cot x = \frac {1}{\tan x}, \quad \sec x = \frac {1}{\cos x}, \quad \operatorname{cosec} x = \frac {1}{\sin x}.
\]

Recall (Gen.~Top., VIII, p.~109) that the angle corresponding to the base \(2\pi\) is called the radian.

